\documentclass[12pt]{article}
\usepackage[utf8]{inputenc}
\usepackage[T1]{fontenc}
\usepackage[margin=1in]{geometry}
\usepackage{setspace}
\usepackage{xcolor}
\usepackage[colorlinks=true, urlcolor=blue!60!black, linkcolor=blue!60!black]{hyperref}
\setlength{\parindent}{0em}
\setlength{\parskip}{0.45em}
\renewcommand{\baselinestretch}{1.0}
\begin{document}
\begin{center}
{\Large \textit{The Civil Rights Cases}, 109 U.S. 3 (1883)\par}
\vspace{0.3em}
{\normalsize Case Brief --- One-Page Edition\par}
\end{center}

\textbf{Memory Jogger.} State action doctrine: Congress could not ban private discrimination in inns, theaters, and railroads (\href{https://www.courtlistener.com/opinion/90897/civil-rights-cases/}{CourtListener}).

\textbf{Facts.} The Civil Rights Act of 1875 guaranteed all persons ``full and equal enjoyment'' of inns, public conveyances, and theaters. Five consolidated cases involved innkeepers in Kansas and Missouri, theaters in San Francisco and New York, and a conductor who barred Sallie Robinson, a Black woman, from a railroad's ladies' car.

\textbf{Procedural History.} In three cases the circuit judges divided and certified the question; in California a demurrer was sustained, and in Tennessee a jury found for the railroad.

\textbf{Votes.} 8--1. Bradley wrote for the Court; Harlan dissented.

\textbf{Holding.} Sections 1 and 2 are void. The Fourteenth Amendment reaches only state action, and Section 5 permits only corrective legislation; private denial of accommodations is not a badge or incident of slavery.

\textbf{Reasoning.} The Act reached private individuals even in States with fair laws, creating a federal ``code of municipal law'' contrary to the Tenth Amendment. The Thirteenth Amendment reaches slavery's inseparable incidents, but treating every private refusal as one would be ``running the slavery argument into the ground.'' The Commerce Clause question was reserved. Harlan called the reasoning ``entirely too narrow and artificial,'' noting that the Court had once implied broad federal power to protect slaveholders, and treating carriers, inns, and licensed theaters as ``agents or instrumentalities of the State.''

\textbf{Examples.} \emph{Same-Side:} a federal remedy against private gender-motivated violence (\textit{Morrison}). \emph{Opposite-Side:} court enforcement of racial covenants (\textit{Shelley}). \emph{Fence-Sitter:} a segregating restaurant in a city garage.

\textbf{Critique.} Defenders cite the text ``No State shall'' and federalism; originalists such as Michael McConnell argue Reconstruction intended more; progressives place the case in the anti-canon, and Charles Black called the public/private line illusory.

\textbf{Key Quotation.} ``It is State action of a particular character that is prohibited. Individual invasion of individual rights is not the subject matter of the amendment.'' --- Justice Bradley, 109 U.S. at 11.

\end{document}
